\section{Conclusion}
\label{sec:conclusion}

We develop \ours, a Bayesian evidence quantification framework for diagnosing null
results in AI alignment research. \ours computes a closed-form Bayes factor
$\mathrm{BF}_{01}$ from observed effect sizes and sample sizes, directly addressing
the fundamental ambiguity between genuine model robustness and experimental
insufficiency.

Applied to 27 published alignment null results, \ours finds that zero studies provide
strong robustness evidence, 85\% are inconclusive, and \mainresult have $P(\text{weakness}) > 0.3$.
The field is systematically producing uninformative null results that researchers
interpret as safety evidence. The four studies with genuine (moderate) robustness
evidence are either structural nulls or geometrically constrained---not typical
empirical demonstrations of robustness.

Prospective validation with real \llama experiments confirms the framework's practical
value through the inverse jailbreak effect: strong explicit jailbreak patterns achieve
0\% success while weak educational framing achieves 20\%. \ours correctly classifies
the strong-intervention null as inconclusive ($P(\text{weakness}) = 0.45$), preventing
a false safety claim that a naive reading of ``0 jailbreaks succeeded'' would license.

{\bf Key takeaway.} The alignment field needs higher sample sizes, more realistic
intervention designs, and principled evidence quantification. The typical alignment
study, at its current scale, cannot demonstrate robustness---it can only fail to
detect vulnerability, which is not the same thing.

{\bf Future work.} The most impactful extensions are: (1) a hierarchical Bayesian
model with intervention-type-specific effect size priors, calibrated from a formal
meta-analysis of $100+$ alignment studies; (2) large-scale prospective experiments
with $n \geq 50$ per condition to validate sequential stopping rules in practice;
and (3) integration with existing evaluation frameworks such as RewardBench
\citep{lambert2024rewardbench} and OpenAI Evals, embedding the \ours diagnostic
as a post-hoc analysis layer.
